\documentclass[12pt,a4paper,fontset=none]{ctexart}

\usepackage{fontspec}
\usepackage{xeCJK}
\setmainfont{Times New Roman}
\setCJKmainfont{SimSun}[BoldFont=SimHei]

\usepackage{geometry}
\geometry{a4paper, margin=2.5cm}

\usepackage{amsmath,amssymb,amsthm}
\usepackage{xcolor}
\usepackage{mdframed}
\usepackage{etoolbox}
\usepackage{fancyhdr}
\usepackage{titlesec}
\usepackage{hyperref}

\usepackage{tocloft}
\renewcommand{\cftsecleader}{\hfill}
\renewcommand{\cftsubsecleader}{\cftdotfill{\cftdotsep}}

\definecolor{capellaBlue}{RGB}{0,102,204}
\definecolor{capellaRed}{RGB}{204,0,0} % new example color

\mdfdefinestyle{theoremstyle}{
    linecolor=capellaBlue,
    leftline=true,
    rightline=false,
    topline=false,
    bottomline=false,
    backgroundcolor=capellaBlue!6,
    linewidth=2pt,
    innertopmargin=6pt,
    innerbottommargin=6pt,
    skipabove=4pt,
    skipbelow=4pt
}

% identical frame style for examples but in red
\mdfdefinestyle{examplestyle}{
    linecolor=capellaRed,
    leftline=true,
    rightline=false,
    topline=false,
    bottomline=false,
    backgroundcolor=capellaRed!6,
    linewidth=2pt,
    innertopmargin=6pt,
    innerbottommargin=6pt,
    skipabove=4pt,
    skipbelow=4pt
}

\newtheoremstyle{bluehead}
  {0pt}{0pt}
  {\normalfont}
  {}
  {\bfseries\color{capellaBlue}}
  {}
  { }
  {\thmname{#1}\ \thmnumber{#2}\ \thmnote{\ {\textcolor{capellaBlue}{(#3)}}}}

% red heading for examples
\newtheoremstyle{redhead}
  {0pt}{0pt}
  {\normalfont}
  {}
  {\bfseries\color{capellaRed}}
  {}
  { }
  {\thmname{#1}\ \thmnumber{#2}\ \thmnote{\ {\textcolor{capellaRed}{(#3)}}}}

\newcounter{theorem_counter}

% definitions in blue
\theoremstyle{bluehead}
\newtheorem{definition}{定义}[subsection]
\surroundwithmdframed[style=theoremstyle]{definition}

% examples in red (same formatting except color)
\theoremstyle{redhead}
\newtheorem{example}{例题}[subsection]
\surroundwithmdframed[style=examplestyle]{example}


\pagestyle{fancy}
\fancyhf{}
\fancyhead[L]{ CHM NOTE}
\fancyhead[R]{\leftmark}
\fancyfoot[C]{\thepage}


\begin{document}

\begin{titlepage}
\centering
\vspace*{\fill}
\begin{minipage}{\textwidth}
  \centering
  {\Huge\bfseries 预科二下半学期数学笔记\par}
  \vspace{1cm}
  {\Large 堡堡菲\par}
  \vspace{1cm}
    {\large \today\par}
\end{minipage}
\vspace*{\fill}
\end{titlepage}


\tableofcontents
\thispagestyle{empty}
\newpage

\section*{导言}
\addcontentsline{toc}{section}{导言}

我说把 设$\varepsilon>0$ 写成 $\forall \varepsilon>0$ 的入平时成绩零分, 你们尔朵龙吗?

\newpage


\section{数列极限}

\subsection{数列极限例题补充}

鉴于 CHM 在本节课中讲解的内容 WLZ 上学期几乎全部涉及, 我就挑选一道例题.

\vspace{0.2cm}

\begin{example}
  设 $|a|>1$, $\alpha \in \mathbb{R^+}$, 那么 $\displaystyle \lim_{n \to \infty}\frac{n^\alpha}{a^n}=0$.
\end{example}

将 $a^n$ 形式上记为 $(1+(a-1))^n$ 并利用二项式定理展开, 保留第 $[\alpha]+1$ 项即可.  

\vspace{0.2cm}
\hspace{-2em} \textbf{习题 A} 证明恒等式: $\displaystyle C_{n}^{k-1}+C_{n}^{k}=C_{n+1}^k$.

\vspace{0.2cm}

\hspace{-2em} \textbf{习题 B} (Cesàro 求和极限)

\vspace{0.2cm}

设 $\{a_n\}_{n \ge 1}$ 为实数序列, 记 $\displaystyle \sigma_n=\frac{a_1+a_2+\cdots +a_n}{n}$.

\vspace{0.2cm}

\hspace{-2em} 1) 证明 $\displaystyle \lim_{n \to \infty} a_n=a \Rightarrow \displaystyle \lim_{n \to \infty} \sigma_n=a$. 其逆命题是否成立?

\hspace{-2em} 2) 对 $k \ge 1$, 记 $b_k=a_{k+1}-a_{k}$. 证明, $\forall n \in \mathbb{N}_{\ge 2}$, 都有 $\displaystyle a_n-\sigma_n=\frac{1}{n}\sum_{k=1}^{n-1}kb_k$.

\hspace{-2em} 3) 设 $\{kb_k\}_{k \ge 1}$ 有界并且 $\{\sigma_{n}\}_{n \ge 1}$ 收敛, 证明 $\{a_n\}$ 也收敛.


\subsection{秘}

\LaTeX 秘.

\newpage

\section{第二章}

\begin{definition}
目录会自动显示本节标题。
\end{definition}

\end{document}